\subsection{连续线性映射的转置}

在本节中，\(E_{1}\) 与 \(E_{2}\) 表示两个局部凸空间，其对偶分别为 \(E_{1}^{\prime}\) 与 \(E_{2}^{\prime}\)。

设 \(u\) 是从 \(\mathrm{E}_1\) 到 \(\mathrm{E}_2\) 中的线性映射。为使 \(u\) 在 \(\mathrm{E}_1\) 与 \(\mathrm{E}_2\) 赋予弱化拓扑后连续，必要且充分的是：\(f \circ u\) 属于 \(\mathrm{E}_1'\)（对每个 \(f \in \mathrm{E}_2'\)）；若 \(u\) 连续，则此条件成立。于是，线性映射 \(f \mapsto f \circ u\)（它把 \(\mathrm{E}_2'\) 映到 \(\mathrm{E}_1'\) 中）称为 \(u\) 的转置，记作 \(^t u\)。

命题 5. --- 设 u 是从 \(E_{1}\) 到 \(E_{2}\) 中的连续线性映射。

\begin{enumerate}
\def\labelenumi{(\roman{enumi})}
\item
  若 \(E_{1}\) 与 \(E_{2}\) 是 Hausdorff 的，则 u 为单射当且仅当 \(^{t}u\) 的像在 \(E_{1}^{\prime}\) 中对于弱拓扑 \(\sigma(\mathrm{E}_{1}^{\prime},\mathrm{E}_{1})\) 稠密。
\item
  为使 \({}^{t}u\) 单射，必要且充分的是 \(u\) 的像在 \(\mathbf{E}_2\) 中稠密。
\end{enumerate}

\(E_{2}\) 的向量子空间对于初始拓扑稠密，当且仅当它对于弱化拓扑稠密（IV，第 4 页，命题 2）。于是命题 5 由 II 第 47 页的推论 2 得出。

命题 6. --- 设 u 是从 \(E_{1}\) 到 \(E_{2}\) 中的线性映射，且它对于弱化拓扑连续。对 i = 1, 2，设 \(S_{i}\) 是 \(E_{i}\) 的有界子集的一个族。为使 \(^{t}u\) 是从 \((\mathrm{E}_{2}^{\prime})_{\mathfrak{S}_{2}}\) 到 \((\mathrm{E}_{1}^{\prime})_{\mathfrak{S}_{1}}\) 中的连续映射，必要且充分的是：对每个集 \(A \in S_{1}\)，存在集 \(A_{1}, \ldots, A_{n}\)（均属于 \(S_{2}\)）以及一个实数 \(\lambda > 0\)，使得 \(\lambda.u(\mathbf{A})\) 包含在 \(A_{1} \cup \ldots \cup A_{n}\) 的闭均衡凸包中。

这是 III 第 15 页命题 2 的直接推论。

推论. --- 设 u 是从 \(E_{1}\) 到 \(E_{2}\) 中的连续线性映射。则 \({}^t u\) 在对偶空间 \(E_{i}'\) 赋以下列拓扑时连续：

\begin{enumerate}
\def\labelenumi{\alph{enumi})}
\item
  弱拓扑 \(\sigma (\mathrm{E}_i^{\prime},\mathrm{E}_i)\)
\item
  强拓扑 \(\beta (\mathrm{E}_i^{\prime},\mathrm{E}_i)\)
\item
  Mackey 拓扑 \(\tau (\mathrm{E}_i^{\prime},\mathrm{E}_i)\)
\item
  预紧收敛拓扑。
\end{enumerate}

此外，若 \(E_{2}\) 是 Hausdorff 的，则 \(^{t}u\) 在对偶空间 \(E_{i}^{\prime}\) 赋以下列拓扑时连续：

\begin{enumerate}
\def\labelenumi{\alph{enumi})}
\setcounter{enumi}{4}
\tightlist
\item
  紧收敛拓扑（相应地，紧凸收敛拓扑）。
\end{enumerate}

唯一需要证明的是情形 \(c\)），即 \(\mathrm{E}_1\) 与 \(\mathrm{E}_2\) 的拓扑未必为 Hausdorff 的情况。此时，对每个线性型 \(f \in \mathrm{E}'_1{}^*\)，\(f \circ {}^t u\) 是 \(\mathrm{E}_2'\) 上的线性型；于是存在线性映射 \(v: \mathrm{E}'_1{}^* \to \mathrm{E}'_2{}^*\)，它对于拓扑 \(\sigma(\mathrm{E}'_1{}^*, \mathrm{E}'_1)\) 与 \(\sigma(\mathrm{E}'_2{}^*, \mathrm{E}'_2)\) 连续，且满足 \(d_{\mathrm{B}_2} \circ u = v \circ d_{\mathrm{B}_1}\)，其中 \(d_{\mathrm{B}_i}\) 是从 \(\mathrm{E}_i\) 到 \(\mathrm{E}'_i{}^* (i = 1, 2)\) 中的典范映射。因此，若 A 是 \(\mathrm{E}_1\) 的子集，且 \(d_{\mathrm{B}_1}(\mathrm{A})\) 对于 \(\sigma(\mathrm{E}'_1{}^*, \mathrm{E}'_1)\) 是凸、均衡且紧的，则 \(d_{\mathrm{B}_2}(u(\mathrm{A})) = v(d_{\mathrm{B}_1}(\mathrm{A}))\) 对于 \(\sigma(\mathrm{E}'_2{}^*, \mathrm{E}'_2)\) 是凸、均衡且紧的，因为拓扑 \(\sigma(\mathrm{E}'_1{}^*, \mathrm{E}'_1)\) 与 \(\sigma(\mathrm{E}'_2{}^*, \mathrm{E}'_2)\) 都是 Hausdorff 的。

命题 7. --- 设 \(u: \mathrm{E}_1 \to \mathrm{E}_2\) 是线性映射。我们假定 \(u\) 对于 \(\mathrm{E}_1\) 与 \(\mathrm{E}_2\) 的弱化拓扑连续。

\begin{enumerate}
\def\labelenumi{(\roman{enumi})}
\item
  映射 \(u\) 在 \(\mathrm{E}_1\) 与 \(\mathrm{E}_2\) 赋予各自的 Mackey 拓扑时是连续的。
\item
  若 \(\mathbf{E}_1\) 是有界型的或桶型的，则 \(u\) 对于 \(\mathbf{E}_1\) 与 \(\mathbf{E}_2\) 的初始拓扑连续。
\item
  为使 \(u\) 对于 \(\mathbf{E}_1\) 与 \(\mathbf{E}_2\) 的初始拓扑连续，必要且充分的是：在 \(^t u\) 之下，\(\mathbf{E}_2'\) 的每个等度连续子集的像在 \(\mathbf{E}_1'\) 中等度连续。
\end{enumerate}

该假设蕴含 \({}^t u\) 对于弱拓扑 \(\sigma(\mathrm{E}_2', \mathrm{E}_2)\) 与 \(\sigma(\mathrm{E}_1', \mathrm{E}_1)\) 连续（II，第 46 页，推论），因此 \({}^t u\) 将对于 \(\sigma(\mathrm{E}_2', \mathrm{E}_2)\) 凸、均衡且紧的子集映为对于 \(\sigma(\mathrm{E}_1', \mathrm{E}_1)\) 凸、均衡且紧的集，这是因为拓扑 \(\sigma(\mathrm{E}_2', \mathrm{E}_2)\) 与 \(\sigma(\mathrm{E}_1', \mathrm{E}_1)\) 都是 Hausdorff 的。于是，断言 (i) 由 GT，X，§ 1，第 4 目，命题 3，b) 得出。断言 (ii) 是 (i) 的推论：因为若 \(\mathrm{E}_1\) 是有界型的或桶型的，则其初始拓扑就是 Mackey 拓扑，而 \(\mathrm{E}_2\) 的 Mackey 拓扑比 \(\mathrm{E}_2\) 的初始拓扑细。最后，\(\mathrm{E}_i\) 的初始拓扑是在 \(\mathrm{E}_i'\) 的等度连续子集上一致收敛的拓扑（III，第 19 页，命题 7 的推论 1）。这就证明了 (iii)。

推论. --- 设 \(E_{1}\) 是赋范空间。设 u 是从 \(E_{1}\) 到 \(E_{2}\) 中的线性映射。下列性质等价：

\begin{enumerate}
\def\labelenumi{\alph{enumi})}
\item
  \(u\) 连续；
\item
  \(u\) 对于弱化拓扑连续；
\item
  \(\mathrm{E}_1\) 中单位球在 \(u\) 下的像在 \(\mathrm{E}_2\) 中有界；
\item
  对每个点列 \((x_{n})\)（由 \(\mathbf{E}_1\) 中对于初始拓扑趋于 0 的点构成），序列 \((u(x_{n}))\) 对于 \(\mathbf{E}_2\) 的弱化拓扑有界。
\end{enumerate}

由于 \(E_{1}\) 是有界型的，a) 与 b) 的等价性由命题 7 得出；a) 与 c) 的等价性是显然的。a) 与 d) 的等价性由 IV 第 1 页的命题 1 以及 III 第 11 页的命题 1 得出。

命题 8. --- (i) 设 E 是赋范空间，其对偶为 E'。对每个 \(x \in E\)，我们有

\[
\| x \| = \sup _ {x ^ {\prime} \in \mathrm{E} ^ {\prime}, \| x ^ {\prime} \| \leqslant 1} | \langle x, x ^ {\prime} \rangle |.\tag{3}
\]

\begin{enumerate}
\def\labelenumi{(\roman{enumi})}
\setcounter{enumi}{1}
\tightlist
\item
  设 \(\mathbf{E}_1\) 与 \(\mathbf{E}_2\) 是两个赋范空间，\(u\) 是从 \(\mathbf{E}_1\) 到 \(\mathbf{E}_2\) 中的连续线性映射。我们有
\end{enumerate}

\[
\left\| ^ {t} u \right\| = \left\| u \right\|.\tag{4}
\]

设 \(x \in E\)。对每个 \(x' \in E'\)（满足 \(\| x' \| \leqslant 1\)），我们有

\[
\left| \langle x, x ^ {\prime} \rangle \right| \leqslant \| x \| \| x ^ {\prime} \| \leqslant \| x \|.
\]

由 Hahn-Banach 定理（II，第 23 页，推论 2），存在元素 \(x'\)（属于 \(E'\)），使得 \(\| x' \| \leqslant 1\) 且 \(\langle x, x' \rangle = \| x \|\)。这就证明了 (i)。

现在证明 (ii)。由公式 (3) 以及转置的定义，我们有

\[
\begin{array}{l} \| ^ {t} u \| = \sup _ {\| y ^ {\prime} \| \leqslant 1} \left\| ^ {t} u (y ^ {\prime}) \right\| = \sup _ {\| y ^ {\prime} \| \leqslant 1, \| x \| \leqslant 1} \left| \langle x, ^ {t} u (y ^ {\prime}) \rangle \right| \\ = \sup _ {\| x \| \leqslant 1, \| y ^ {\prime} \| \leqslant 1} \left| \langle u (x), y ^ {\prime} \rangle \right| = \sup _ {\| x \| \leqslant 1} \left\| u (x) \right\| = \| u \|. \end{array}
\]

注. --- 1) 公式 (3) 是 (4) 的特殊情形，对应于从 K 到 E 中的线性映射 \(\lambda \mapsto \lambda x\)。

\begin{enumerate}
\def\labelenumi{\arabic{enumi})}
\setcounter{enumi}{1}
\tightlist
\item
  令 \(\mathrm{B}(x, y') = \langle u(x), y' \rangle = \langle x, {}^t u(y') \rangle\)（其中 \(x \in \mathrm{E}_1\)，\(y' \in \mathrm{E}_2'\)）。上述证明表明 \(\mathbf{B}\) 是 \(\mathrm{E}_1 \times \mathrm{E}_2'\) 上的连续双线性型，其范数（GT，X，§ 3，第 2 目）等于 \(\| u \|\)。
\end{enumerate}

推论. --- 设 E 是满足第一可数公理的赋范空间。存在 \(E' - \{0\}\) 的可数子集 D，使得我们有

\[
\| x \| = \sup _ {\xi \in \mathbf {D}} | \langle x, \xi \rangle | / \| \xi \|\tag{5}
\]

对所有 \(x \in E\) 成立。

设 \(B'\) 是对偶 \(E'\)（即 \(E\) 的对偶）的单位球，并对其赋予弱拓扑 \(\sigma(E', E)\)。则 \(B'\) 是紧度量空间（III，第 19 页，推论 2）；因此存在可数稠密子集 \(D'\)（在 \(B'\) 中）。令 \(D = D' \cap (E' - \{0\})\)。设 \(x \in E\)；映射 \(x' \mapsto \langle x, x' \rangle\)（它把 \(B'\) 映到 \(K\) 中）连续，因此

\[
\sup _ {x ^ {\prime} \in \mathbf {B} ^ {\prime}} | \langle x, x ^ {\prime} \rangle | = \sup _ {\xi \in \mathbf {D} ^ {\prime}} | \langle x, \xi \rangle | \leqslant \sup _ {\xi \in \mathbf {D}} | \langle x, \xi \rangle | / \| \xi \| \leqslant \| x \|.
\]

于是公式 (5) 由 (3) 得出。

